\begin{onehalfspace}

Il rumore può essere affrontato in momenti diversi della pipeline. Distinguere questi livelli evita due equivoci: una stima a posteriori dell'output ideale non corregge lo stato durante il calcolo, e un modello appreso non coincide con la correzione d'errore quantistica.

\section{Panoramica e criteri di scelta}

\begin{table}[!htbp]
\centering
\small
\begin{tabular}{@{}p{2.8cm}p{2.8cm}p{3.6cm}p{3.6cm}@{}}
\toprule
\textbf{Strategia} & \textbf{Livello} & \textbf{Principio} & \textbf{Vincolo principale} \\
\midrule
Soppressione del rumore
& hardware e controllo
& ridurre l'errore prima o durante l'esecuzione
& richiede accesso a pulse, calibrazione o dispositivo \\[2pt]
Mitigazione
& esecuzioni e post-processing
& stimare l'osservabile ideale da circuiti rumorosi
& aumenta il numero di campioni e dipende dal modello \\[2pt]
\ac{QEC} e fault tolerance
& stato logico e architettura
& rilevare errori tramite sindromi e limitarne la propagazione
& richiede qubit, cicli di sindrome e operazioni logiche aggiuntive \\[2pt]
Metodi appresi classici
& output o decoder
& apprendere una decisione da istogrammi o sindromi
& necessita dati, validazione e un riferimento non appreso \\
\bottomrule
\end{tabular}
\caption{Livelli di intervento contro il rumore e relativi vincoli.}
\label{tab:strategie_confronto}
\end{table}

La \textbf{soppressione} comprende isolamento, calibrazione, pulse shaping e dynamical decoupling. Riduce l'errore fisico ma non lo annulla ed è in larga parte vincolata all'hardware. La \textbf{mitigazione}, fra cui la \ac{ZNE}, combina più esecuzioni rumorose senza codificare qubit logici; è utilizzabile nel regime NISQ, ma il suo costo di campionamento cresce con il rumore totale e con la tecnica adottata.

La \textbf{\ac{QEC}} codifica l'informazione in un sottospazio logico e ricava una sindrome senza misurare direttamente lo stato protetto. La fault tolerance estende questo principio all'intero calcolo, imponendo che un singolo guasto non si propaghi in modo incontrollato. Sotto le ipotesi del threshold theorem e al di sotto della soglia pertinente, aumentare la distanza del codice può sopprimere l'errore logico~\cite{aharonov1997fault}; soglia e overhead dipendono però da circuito, modello di rumore e decoder.

I \textbf{metodi appresi} sono trasversali. Possono filtrare un output già campionato oppure assistere il decoder di un codice, ma non rimuovono il costo di generare quei dati e non sostituiscono automaticamente una regola analitica. La loro utilità deve essere valutata sullo stesso input e con la stessa metrica dell'alternativa di confronto.

\section{Machine learning classico su dati quantistici}
\label{sec:qml_overview}

In questa tesi Random Forest, \ac{SVM} e \ac{MLP} sono modelli eseguiti su hardware classico; il termine corretto è quindi \textbf{machine learning classico applicato a dati quantistici}, non \ac{QML}. Vengono considerati due contratti distinti:
\begin{itemize}
    \item nel \textbf{post-processing di Shor}, il modello riceve feature dell'istogramma e decide se ampliare la ricerca dal candidato dominante ai primi $K$ candidati;
    \item nella \textbf{decodifica QEC}, il modello riceve sindromi o detection event e propone una correzione, oppure modifica residualmente l'uscita di un decoder analitico.
\end{itemize}
Nel primo caso il riferimento naturale è una regola TOP-$K$ priva di addestramento; nel secondo è un decoder analitico adatto al codice e al noise model. Separare i due contratti impedisce di attribuire al modello informazione già contenuta nella regola di selezione o nel decoder di base.

\section{Correzione quantistica e decodifica}
\label{sec:qec_overview}

La \ac{QEC} distribuisce un qubit logico su più qubit fisici. Gli stabilizzatori rivelano la classe dell'errore attraverso una sindrome, mentre il decoder sceglie un recupero compatibile. Questo è sufficiente per motivare gli esperimenti successivi; formalismo, circuiti e curve dei codici a ripetizione e di Steane sono sviluppati nel Capitolo~\ref{chap:qec}, mentre geometria, soglia e distanza del surface code sono trattate nel Capitolo~\ref{chap:surface}. I riferimenti storici dei codici impiegati sono Shor~\cite{shorcorrection1995}, Steane~\cite{steane1996} e il surface code~\cite{fowler2012}.

\subsection{Il modello appreso come componente del decoder}
\label{sec:qec_ml}

Il punto di contatto fra \ac{QEC} e machine learning è la mappa classica dalla sindrome alla decisione. Un modello può apprendere correlazioni non rappresentate dal decoder analitico oppure operare come correzione residua~\cite{torlaiMelko2017,alphaqubit2024}. Ciò non garantisce un vantaggio: la qualità dipende dalla distribuzione di addestramento, dalla stabilità del rumore e dal costo d'inferenza. Per questo il Capitolo~\ref{chap:surface} confronta decoder appreso e analitico sugli stessi campioni e mantiene separati addestramento, selezione e valutazione.

\section{Strategia sperimentale della tesi}
\label{sec:motivazione_ml}

Il lavoro adotta una progressione in due fasi. La prima studia ciò che si può recuperare \emph{dopo} la misura di Shor: TOP-1, TOP-$K$ e filtro appreso sono confrontati sugli stessi istogrammi, mentre la \ac{ZNE} costituisce un confronto di mitigazione separato. Metodo e verifica sono descritti nei Capitoli~\ref{chap:metodologia} e~\ref{chap:risultati}.

La seconda fase interviene \emph{durante} la protezione dell'informazione: caratterizza codici di complessità crescente, stima l'errore logico e valuta il ruolo di un decoder appreso. Il Capitolo~\ref{chap:integrazione} mantiene infine distinta la sensibilità di Shor a un proxy per gate dai tassi logici per ciclo dei codici. Questa separazione fra livelli --- output, sindrome e circuito logico --- è il criterio che rende confrontabili le strategie senza confonderne obiettivi e metriche.

\end{onehalfspace}
